\chapter{Positive measure, empty interior and infinite perimeter}\label{sec:fat-cantor-examples}
A positive-length Cantor set separates measure from topology and rectifiability from finite boundary mass. The construction is in \cite{Mor16}, fifth edition, Example after Proposition 3.12, p.~33; the associated current is Exercise 4.17, p.~59. We give the construction and the variation argument explicitly.

\begin{definition}[A fat Cantor construction]\label{def:fat-cantor-set}
Start with $E_0=[0,1]$. At step $n\ge1$, remove the central open interval of length $4^{-n}$ from each of the $2^{n-1}$ closed components of $E_{n-1}$. Let $E_n$ be the remaining union of $2^n$ closed intervals and define $E=\bigcap_{n\ge0}E_n$.
\end{definition}

\begin{proposition}[Length and topology of the fat Cantor set]\label{prop:fat-cantor-length-topology}
The construction in Definition \ref{def:fat-cantor-set} is well defined. Every stage-$n$ component has length
\[
 \ell_n=2^{-n-1}+2^{-2n-1},
\]
and its intersection with $E$ has length $2^{-n-1}$. The set $E$ is compact, perfect, nowhere dense, and has Lebesgue measure $1/2$.
\end{proposition}
\begin{proof}
The recurrence $\ell_n=(\ell_{n-1}-4^{-n})/2$, with $\ell_0=1$, gives the displayed formula by induction. It is positive and the removed length is smaller than its parent interval, validating every step. At step $n$ the removed total length is $2^{n-1}4^{-n}=2^{-n-1}$. These intervals are disjoint, so their total length is $1/2$, proving $|E|=1/2$ by continuity from above. Within a fixed stage-$n$ interval, the total future removed length is
\[
 \sum_{j\ge1}2^{j-1}4^{-(n+j)}=\tfrac12 4^{-n}.
\]
Subtracting it from $\ell_n$ gives $2^{-n-1}$. Compactness follows by nested closed intersections. Since $\ell_n\to0$, no interval can be contained in every $E_n$, so $E$ has empty interior and, being closed, is nowhere dense. Endpoints of each component survive every later step. For any $x\in E$, choose a stage-$n$ component containing it and an endpoint different from $x$; that endpoint belongs to $E$ and is within $\ell_n$. Hence $E$ has no isolated points.
\end{proof}

\begin{example}[Ordinary continuity and approximate continuity differ]\label{ex:fat-cantor-approximate-continuity}
The function $\chi_E$ is discontinuous at every point of $E$ and continuous at every point outside $E$. Nevertheless its approximate value is one at almost every point of $E$ and zero at almost every point outside $E$. Thus a set of positive measure can consist entirely of ordinary discontinuities of a bounded measurable function that is approximately continuous almost everywhere.
\end{example}
\begin{proof}
Every neighborhood of a point of $E$ meets its dense open complement, where the function is zero, while its value on $E$ is one. Off the closed set $E$, a small neighborhood lies in its complement and the function is constant. The approximate-value statement is the Lebesgue density theorem applied to $E$ and its complement. More explicitly, at a density-one point of $E$, for every $0<\varepsilon<1$, the set $\{|\chi_E-1|>\varepsilon\}$ has density zero; this is the definition of approximate convergence to one. The analogous assertion gives zero off $E$ almost everywhere. No continuity representative is claimed.
\end{proof}

\begin{proposition}[The characteristic function has infinite variation]\label{prop:fat-cantor-infinite-variation}
The function $\chi_E$, extended by zero to $\mathbb R$, does not belong to $BV(\mathbb R)$; equivalently $E$ has infinite one-dimensional perimeter.
\end{proposition}
\begin{proof}
Fix a positive integer $N$. Choose a stage with at least $N+1$ retained components, and select $N+1$ successive such components. Each has positive $E$ measure by Proposition \ref{prop:fat-cantor-length-topology}, so it contains a density-one point $x_i$ of $E$, $0\le i\le N$. Between successive components choose a point $y_i$ in a deleted open gap. In increasing order the selected points are $x_0<y_1<x_1<\cdots<y_N<x_N$.

Let $u_\varepsilon=\rho_\varepsilon*\chi_E$ for a nonnegative smooth compactly supported approximate identity. At the chosen density points $u_\varepsilon(x_i)\to1$, and at the gap points $u_\varepsilon(y_i)\to0$. The fundamental theorem of calculus on each consecutive interval implies
\[
 \int_{\mathbb R}|u_\varepsilon'|\ge
 \sum_{i=1}^N\bigl(|u_\varepsilon(y_i)-u_\varepsilon(x_{i-1})|
                  +|u_\varepsilon(x_i)-u_\varepsilon(y_i)|\bigr)\longrightarrow2N.
\]
If $\chi_E$ were BV, Lemma \ref{lem:bv-mollification} would bound the left side by $|D\chi_E|(\mathbb R)$ for all $\varepsilon$. Since $N$ is arbitrary this is impossible. The equivalence with perimeter is the definition of the perimeter of a measurable set via its characteristic function.
\end{proof}

\begin{example}[A finite-mass rectifiable current with infinite boundary mass]\label{ex:fat-cantor-nonnormal-current}
Define a one-dimensional current on $\mathbb R$ by $T(f\,dx)=\int_E f(x)\,dx$. It is an integer-multiplicity rectifiable current of mass $1/2$, but its boundary has infinite mass. Consequently it is neither normal nor integral.
\end{example}
\begin{proof}
The set $E$ is a measurable subset of the oriented line, with multiplicity one and length $1/2$, giving the rectifiable representation and mass. For a smooth compactly supported function $\varphi$, $\partial T(\varphi)=T(d\varphi)=\int_E\varphi'$. If this zero-dimensional boundary current had finite mass, then
\[
 \left|\int_E\varphi'\right|\le C\|\varphi\|_\infty
\]
for one finite constant $C$ and every such test. This is precisely bounded distributional variation of $\chi_E$, contradicting Proposition \ref{prop:fat-cantor-infinite-variation}. A normal current must have both finite mass and finite boundary mass; an integral current has locally integral boundary and, here with compact support, finite boundary mass. Thus neither property holds. This resolves the distinction posed in \cite{Mor16}, Exercise 4.17.
\end{proof}
